% conclusion.tex -- Section 14: Conclusion and Future Work

\section{Conclusion}
\label{sec:conclusion}

We have introduced Kant as an overall framework for both training and evaluation of ethical reasoning using language models. Key changes made: - "have introduced" to "introduced" - "
framework makes the following core contributions:

\begin{enumerate}[nosep]
    \item \textbf{Comprehensive game library}: 90+ games spanning nine
Strategic areas ranging from classic dilemmas via competition in markets to auction mechanisms and cooperative game theory offer comprehensive coverage of the terrain of strategic reasoning.
    \item \textbf{Multi-agent infrastructure}: $N$-player game support,
A system for forming coalitions including mechanisms for negotiation, enforcement methods and side payments; alongside a high level governance engine that supports democracy.
          rule modification during play.
    \item \textbf{Training pipeline}: GRPO and DPO training on game
Trajectories for curriculum learning through different game contexts; linking direct alignment of theoretical environments to games. To connect game theory environments directly with
          training.
    \item \textbf{Safety transfer evaluation}: systematic evaluation of
whether training using games theory also enhances performance on external safety benchmarks such as HarmBench and ETHICS and Truthful QA Par
          XSTest, and MT-Bench.
    \item \textbf{Alignment-oriented evaluation}: six normalized metrics
Capturing cooperation, resistance to exploitation, Pareto efficiency, fairness, adaptability as well as strategic reasoning; stratification also involved
          game splitting for robust evaluation.
\end{enumerate}

The framework integrates with the OpenEnv platform, providing a
Gymnasium-compatible API, WebSocket-based communication, and automated
tournament infrastructure. The meta-governance engine---which allows agents
to propose, vote on, and apply rule modifications during play---represents
a novel connection between game-theoretic multi-agent systems and
Constitutional AI~\citep{bai2022constitutional}.

\subsection{Future Work}

\paragraph{Scaling experiments.}
Comparing different stages of a training pipeline at varying levels of models for understanding how game theoretic alignment training relates to model capacity.

\paragraph{Multi-model tournaments.}
Competitions among agents that pit multiple Large Language Models (LLMs) against each other directly allow for comparison of reasoning strategies and ethics. You are an
across model families.

\paragraph{Governance emergence studies.}
Detailed investigation into how governance structures develop when actors are endowed with comprehensive meta governance capability through broad interactions.
periods.

\paragraph{Real-world transfer.}
Assessing if training using games theory enhances performance of agents in realistic scenarios involving multiple agents above formal standards. To assess whether training using game theory

\paragraph{Dynamic curriculum optimization.}
Automated curriculum design that identifies which game domains
contribute most to alignment transfer and adjusts training distribution
accordingly.
